% ARCHIVED AUDITED DRAFT: independently checked by root, 2026-11-19 15:20 UTC.
% Incorporated into part_B.tex, whose version is authoritative.
\begin{saresult}{B.055}{Proved}{Localization and completion with arbitrary coefficient extensions}
Let $R$ be a Noetherian commutative ring and $R\subset S$
an arbitrary unital inclusion. No torsion-free hypothesis is imposed.
For every multiplicative subset $U\subset R$ and every $n$,
\[
 U^{-1}\Omega_n(R,S)
       =\Omega_n(U^{-1}R,U^{-1}S).                         \tag{B55.1}
\]
If $(R,\mathfrak m)$ is local, set $A=\widehat R$ and
$S_A=S\otimes_R A$. Then
\[
 \Omega_n(R,S)\otimes_R A=\Omega_n(A,S_A).                \tag{B55.2}
\]
When the residue field is finite, every coefficient of an
integer-valued polynomial lies in the ring
\[
 C(R,S)=\bigcup_{j\ge0}(R:_S\mathfrak m^j),\qquad
 (R:_S\mathfrak m^j)=\{s\in S:\mathfrak m^js\subset R\}.
                                                               \tag{B55.3}
\]
These rings commute with completion, and $C(R,S)/R$ is
elementwise killed by powers of $\mathfrak m$.
All tensor maps and the existence conditions commute with
the identifications above.
\end{saresult}
\begin{proof}
The degree induction in B.033 proves the forward localization
inclusion over arbitrary rings: the identity
$s^dF(X)-F(sX)$ cancels the terms of highest total degree,
and $s$ is invertible after localization. For the converse,
clear the finitely many coefficient denominators to obtain
$F\in S[X_1,\ldots,X_n]$ taking $R^n$ into $U^{-1}R$.
If $s_1,\ldots,s_N$ are its coefficients, the module
$M=R+Rs_1+\cdots+Rs_N\subset S$ is finite over $R$.
The classes of all values $F(r)$ generate a submodule of
the Noetherian module $M/R$. That submodule is finitely
generated and vanishes after localization. One element
$u\in U$ therefore kills all its generators, and
$uF(R^n)\subset R$. This proves (B55.1), without rational
coefficient denominators or any torsion assumption.

For completion, infinite residue field gives ordinary
polynomial rings by the unit Vandermonde argument. Assume
the residue field is finite. Write $C_j=(R:_S\mathfrak m^j)$.
These are submodules containing $R$, and
$C_jC_k\subset C_{j+k}$, so their union is a ring.
Choose generators $a_1,\ldots,a_s$ of $\mathfrak m^j$.
Flatness of completion applied to
\[
 C_j=\ker\bigl(S\longrightarrow(S/R)^s,
                   x\longmapsto(a_ix\bmod R)_i\bigr)
\]
identifies $C_j\otimes_R A$ with
$(A:_{S_A}(\mathfrak m A)^j)$.

For any coefficient $x$ of $F\in\Omega_n(R,S)$, and any
nonmaximal prime $\mathfrak p$ of $R$, the domain
$R/\mathfrak p$ is infinite. A rectangular grid of
distinct residues gives a Vandermonde determinant
$v\notin\mathfrak p$, and Cramer's rule gives $vx\in R$.
Thus the ideal $\{r\in R:rx\in R\}$ is either $R$ or
has radical $\mathfrak m$. Noetherianity then puts $x$
in some $C_j$. This argument also works over $A$:
a nonmaximal prime of $A$ cannot contract to
$\mathfrak m$, so the same $R$-grids give the required
determinants. It does not require $R$ or $A$ to be a domain.
If there are no nonmaximal primes, the argument says that
every proper denominator ideal has radical $\mathfrak m$,
which is exactly what is needed.

For fixed $j$ and a finite rectangular degree bound, the
integer-valued polynomials with coefficients in $C_j$
are the kernel of finitely many evaluations modulo $R$,
at representatives of the finite ring $R/\mathfrak m^j$.
Changing an input by $\mathfrak m^j$ changes every
monomial value by that ideal, and hence changes the
polynomial value by an element of $R$. Since
$R/\mathfrak m^j=A/(\mathfrak m A)^j$, the same finite
tests apply after completion. Tensor the finite evaluation
kernel, then take the filtered union over $j$ and the
degree bounds. This proves (B55.2) and (B55.3), including
the assertion about completion of the coefficient ring.
Polynomial multiplication and substitution give the stated
compatibility; existence for an $A$-linear combination
follows by the same finite linear combination of tensor
expressions. Faithful flatness detects failure of
injectivity and of the requisite image conditions.
\end{proof}

\begin{saresult}{B.056}{Proved}{Complete local epimorphic algebras supported at one point}
Let $(A,\mathfrak m)$ be a complete Noetherian local ring
with finite residue field. Let $A\subset C$ be a unital
ring epimorphism, and suppose every element of $C/A$ is
killed by a power of $\mathfrak m$. Define
\[
 D=\bigcap_{j\ge0}\mathfrak m^j C,\qquad I=A\cap D.
\]
Then
\[
 C=A+D,\qquad D=\mathfrak mD=D^2=eC
                                                               \tag{B56.1}
\]
for an idempotent $e\in C$, and there is an isomorphism
of $A$-algebras
\[
 C\simeq D\times A/I,\qquad
 a\longmapsto(ea,a+I).                                    \tag{B56.2}
\]
The ideal $D$ is finitely generated over $C$. No finiteness
assumption on $C$ itself is needed.
\end{saresult}
\begin{proof}
First, every map $A\to C/\mathfrak m^jC$ is surjective.
It is a ring epimorphism from the finite local ring
$A/\mathfrak m^j$. After reduction modulo its nilpotent
maximal ideal, an epimorphic algebra over the residue
field is either zero or that field: the linear-functional
argument of B.013 applied to $b\otimes1=1\otimes b$
proves this even for an infinite-dimensional algebra.
Thus $C/\mathfrak m^jC$ equals the image of $A$ plus
its maximal-ideal multiple. Iterating this equality
until the maximal ideal vanishes proves surjectivity.

Fix $c\in C$. For each $j$ the set
\[
 X_j=\{a\in A:c-a\in\mathfrak m^jC\}
\]
is nonempty and is a coset of an ideal containing
$\mathfrak m^j$. These are nested closed subsets of
the compact ring $A=\varprojlim A/\mathfrak m^j$.
For an elementary proof of the needed compactness,
choose $a_j\in X_j$, successively pass to infinite
subsequences with constant residue modulo
$\mathfrak m,\mathfrak m^2,\ldots$, and use completeness
on the resulting compatible residues. Their limit
belongs to every $X_j$. Hence $c-a\in D$ for some
$a\in A$, proving $C=A+D$.

The epimorphism condition gives
$C\otimes_A(C/A)=0$. Since $D/I=C/A$, right exactness
shows that $C\otimes_A I\to C\otimes_A D$ is surjective.
Compose with multiplication onto $D$ to get $D=CI$.
The ideal $I$ is finitely generated over $A$, so $D$
is a finitely generated ideal of $C$.

We next prove $D=\mathfrak mD$. There is an integer
$N\ge1$ for which
\[
                       I\cap\mathfrak m^N A\subset\mathfrak mI.
                                                               \tag{B56.3}
\]
Indeed the descending subspaces
$((I\cap\mathfrak m^jA)+\mathfrak mI)/\mathfrak mI$
of the finite-dimensional vector space $I/\mathfrak mI$
stabilize. Their intersection is zero by Krull
intersection on the finite module $A/\mathfrak mI$,
so their stable value is zero.
For $i\in I\subset D\subset\mathfrak m^NC$, write
$i=a+d$ with $a\in\mathfrak m^NA$ and
$d\in\mathfrak m^ND$, using $C=A+D$.
Then $a\in I\cap\mathfrak m^NA\subset\mathfrak mI$,
so $i\in\mathfrak mD$. Since $D=CI$, this gives
$D=\mathfrak mD$.

Also $D=CI=I+DI$. The quotient $D/DI$ is consequently
a finite $A$-module, generated by the image of $I$,
and it equals its $\mathfrak m$-multiple. Nakayama
gives $D=DI$, hence $D=D^2$.
The determinant argument of B.052 applies to this
finitely generated idempotent ideal and gives $D=eC$.
Finally $(1-e)C=(1-e)A$ because $C=A+D$; the kernel
of $A\to(1-e)A$ is precisely $I$. Orthogonal
idempotents now give (B56.2).
\end{proof}

\begin{saresult}{B.057}{Proved}{The torsion obstruction over a completed domain}
Let $(R,\mathfrak m_R)$ be a nonfield Noetherian local
domain with finite residue field, let $K$ be its fraction
field, and put $A=\widehat R$, $Q=K\otimes_R A$.
Let $A\subset C$ be any unital extension such that
$C/A$ is elementwise killed by powers of the maximal
ideal of $A$. The polynomial coaddition and
comultiplication tensors for $\Omega(A,C)$ are uniquely
determined if and only if
\[
 C=A\quad\text{or}\quad
 \bigl(\dim R=1,\ A\text{ is a domain},\ C=Q\bigr),       \tag{B57.1}
\]
where the second identification is as $A$-algebras.
In both positive cases all multivariable tensor maps
are isomorphisms. In particular no nonzero effective
base torsion survives the uniqueness and existence
requirements.
\end{saresult}
\begin{proof}
For each $c\in C$, choose $j$ with
$\mathfrak m_A^jc\subset A$. A monic null polynomial
for the finite ring $A/\mathfrak m_A^j$ gives $P,cP$
in $\Omega(A,C)$. If the tensor map is injective,
the zero-image tensor $cP\otimes P-P\otimes cP$
and top-coefficient extraction give
$c\otimes1=1\otimes c$ in $C\otimes_A C$.
Thus $A\to C$ is a ring epimorphism, and B.056 applies.
Use its notation $C=A+D$, $D=eC$, $I=A\cap D$.

Set $W=R\setminus\{0\}$. A nonzero element of
$\mathfrak m_R$ kills every class of $C/A$ after
localization, so $W^{-1}C=W^{-1}A=Q$.
Let $T$ be the kernel of $C\to Q$, and let $B$ be
its image. Each element $z\in T$ is killed by a
power of $\mathfrak m_A$: some such power sends
$z$ into $A$, where $W$ acts without torsion by
flatness of completion. Since $D=\mathfrak m_A^jD$
for all $j$, one has $TD=0$. It follows that
$D\cap T=0$, because $e$ acts as the identity on
$D$. We may therefore identify $D$ with its image
in $Q$.

Choose the complete coefficient discrete valuation
ring $V\subset A$ used in B.053: in characteristic
$p$ take $k[[t]]$ for $0\ne t\in\mathfrak m_R$,
and in mixed characteristic take the Cohen ring
with uniformizer $p$. Write $F=\operatorname{Frac}(V)$.
The map $A\to B$ is epic, as a composite of epic
maps. B.052 gives its canonical finite-dimensional
$F$-vector ideal $D'\subset Q$, with
$B=A+D'$ and $B/D'$ finite over $A$.
The image of $D$ in $B/D'$ is a finite $A$-module
and equals its maximal-ideal multiple, so it is
zero by Nakayama. Conversely $B/D$ is a quotient
of $A$, hence finite over $A$. The image there
of the $t$-divisible module $D'$ lies in every
$t^j(B/D)$, hence is zero by Krull intersection.
Consequently $D=D'$ is an ideal of $Q$, finite
dimensional over $F$, and contained in $A[1/t]$.

Let $J=A\cap D$ now mean the intersection of the
images inside $Q$, to distinguish it from $I$,
which was the intersection inside $C$. Then
$I\subset J$, and
\[
                              J/I\simeq T.                \tag{B57.2}
\]
Explicitly, for $a\in J$, subtract its unique lift
in $D\subset C$ from $a\in A\subset C$; this
gives every element of $T$, with kernel $I$.
B.052 says $J$ is a finite $V$-lattice in $D$.
Its quotient $J/I$ is torsion, hence killed by
one power of $t$. Therefore $I$ is also a full
finite $V$-lattice in $D$.

Suppose $D\ne0$ and select a local Artinian factor
$Q_0=e_0Q$ of the finite-dimensional algebra $D$.
The order $A_0=e_0A$ is a complete local quotient
of $A$, and $A_0[1/t]=Q_0$, since
$Q_0\subset A[1/t]$. B.048 and $t$-adic
separatedness show that $A_0$ is finite over $V$,
exactly as in B.053.

If $Q_0$ is nonreduced, use the compatible
coefficient field and finite order $O\subset Q_0$
constructed in B.045. The corrected integrality
argument in B.053 proves that $O$ is finite over
the original $V$: the normalized coefficient DVR
is finite integral over the newly chosen DVR
inside $A_0$, with compatible scalar embeddings.
As $I$ is a full lattice in $D$, some $t^N e_0O$
is contained in $I$. If $u$ is the new uniformizer,
finite generation of $O$ and invertibility of
$t$ on $Q_0$ give
\[
                              u^ce_0O\subset I             \tag{B57.3}
\]
for some $c$. The construction of B.045 now gives
its two polynomials supported on $Q_0\subset D$
with all their values in $I\subset A$.
They belong to $\Omega(A,C)$. Its balanced
derivative functional still applies: arbitrary
polynomials in $\Omega(A,C)$ have their values
in $A$, and projection to $Q_0$ puts those values
in $A_0\subset O$. The input and valuation
bounds use only (B57.3). The same nonzero tensor
kernel is therefore obtained. Thus every selected
factor must be a field.

Assume next that $I$ is a proper ideal of $A$.
For a selected field factor, let $V_1$ be the
finite discrete valuation normalization of $A_0$,
with uniformizer $u$. Clearing its finite lattice
into $I$ gives, after increasing $c$ if necessary,
\[
                     u^ce_0V_1\subset I,\qquad c\ge1.
\]
Put $x=u^ce_0\in I$. By (B56.2), the second
component of any $H\in\Omega(A,C)$ is a
polynomial over $A/I$. Its evaluations at $x$
and $0$ are equal. Since the actual values belong
to $A$, this proves
\[
                         H(x)-H(0)\in I.                 \tag{B57.4}
\]
Set $r=c+1$. An integer-valued congruence indicator
on $V_1$, as constructed in B.024, distinguishes
$0$ and $u^{cr}$. Thus there is a polynomial
$G$ supported on $Q_0$ with all values in
$u^ce_0V_1\subset I$, with $G(0)=0$, and with
selected valuation $v(G(x^r))=c$.
If every binary multiplicative expression existed,
iteration would express $G(X_1\cdots X_r)$ as
a sum of products of one-variable elements of
$\Omega(A,C)$. Taking the difference at $x$ and
$0$ in each variable and using (B57.4) would
put $G(x^r)$ in $I^r$. But $I\subset\mathfrak m_A$,
the map $A\to A_0$ is local, and $V_1$ dominates
$A_0$. Every element of $I^r$ has selected
valuation at least $r=c+1$, a contradiction.
Therefore $D\ne0$ forces $I=A$.

If $D=0$, then $C=A$. If $I=A$, then $1\in D$,
so $C=D$ and its image in $Q$ is the whole of
$Q$. Thus $C=Q$, finite-dimensional over $F$.
B.048 applied to the separated submodule
$A\subset Q=A[1/t]$ makes $A$ finite over $V$;
therefore $\dim R=\dim A=1$.
Nonreduced factors have already been excluded.
If $Q$ had two field factors, the supported
finite-difference argument of B.050, with the
proper ideal $A\cap e_0Q$, would exclude
multiplicative existence. Hence $Q$ is a field,
and its subring $A$ is a domain. This proves
the necessity of (B57.1).

Conversely $C=A$ gives the polynomial ring. In
the second case $Q$ is the fraction field of
the complete one-dimensional domain $A$, whose
normalization is finite and local. B.024 and
B.050 give every multivariable tensor isomorphism
and the required compatible co-operations.
\end{proof}

\begin{saresult}{B.058}{Proved}{Every Noetherian domain with arbitrary coefficient extensions}
Let $R$ be a nonfield Noetherian domain with fraction
field $K$, and let $R\subset S$ be an arbitrary unital
inclusion, allowing base torsion in $S$. For each maximal
ideal $\mathfrak m$ with finite residue field define
the actual coefficient subring of $S_{\mathfrak m}$
\[
 C_{\mathfrak m}
   =\{s\in S_{\mathfrak m}:\mathfrak m^jR_{\mathfrak m}s
              \subset R_{\mathfrak m}\text{ for some }j\}.
\]
The polynomial coaddition and comultiplication tensors
are uniquely determined if and only if, for every such
$\mathfrak m$,
\[
\begin{gathered}
 C_{\mathfrak m}=R_{\mathfrak m},\quad\text{or}\\
 \dim R_{\mathfrak m}=1,\qquad
 \widehat{R_{\mathfrak m}}\text{ is a domain},\qquad
 C_{\mathfrak m}=K.                                      \tag{B58.1}
\end{gathered}
\]
The last equality means the canonical copy of $K$
extending $R_{\mathfrak m}$ inside $S_{\mathfrak m}$.
No torsion-free or finite-normalization hypothesis is
imposed. In every positive case all multivariable
tensor maps are isomorphisms and the compatible
plethory is unique. Together with B.013--B.014, this
settles every Noetherian domain base and every
commutative coefficient extension.
\end{saresult}
\begin{proof}
B.055 permits localization and completion for this
arbitrary extension. At an infinite-residue maximal
ideal the unit Vandermonde argument gives the ordinary
polynomial ring. At a finite-residue maximal ideal,
replace $S_{\mathfrak m}$ by $C_{\mathfrak m}$, which
contains every possible coefficient by B.055.
After completion the latter ring has quotient over
the base killed elementwise by maximal-ideal powers,
so B.057 gives exactly its two positive possibilities.

Faithful flatness descends them. If
$C_{\mathfrak m}\otimes A=A$, then
$(C_{\mathfrak m}/R_{\mathfrak m})\otimes A=0$,
so $C_{\mathfrak m}=R_{\mathfrak m}$.
If $C_{\mathfrak m}\otimes A=K\otimes A$, the
kernel of $C_{\mathfrak m}\to K\otimes C_{\mathfrak m}$
vanishes after faithful flat base change, so it was
zero. Its image lies in $K$, since every coefficient
is sent into $R_{\mathfrak m}$ by a nonzero element
of a maximal-ideal power. Now
$(K/C_{\mathfrak m})\otimes A=0$ proves
$C_{\mathfrak m}=K$ as the canonical subring.
Thus the completed conditions descend precisely
to (B58.1), and conversely these local conditions
give the positive cases of B.057.

The local positive tensor isomorphisms imply the
global ones by testing kernels and cokernels at
every maximal ideal. The existence identities
and their compatibility follow as in B.009.
In the torsion-free case
$C_{\mathfrak m}=S_{\mathrm{rat},\mathfrak m}
\cap R_{\mathfrak m}^{\sharp}$, so this criterion
recovers B.054 while allowing arbitrary additional
torsion that never occurs in an integer-valued
polynomial coefficient.
\end{proof}
